\chapter{Methodology}
\label{chap:methodology}

\section{Problem Formulation}

Given a source language $S$ with labeled dataset
$\mathcal{D}_S = \{(x_i^S, y_i^S)\}_{i=1}^{N_S}$ and a target
language $T$ with limited labeled dataset
$\mathcal{D}_T = \{(x_j^T, y_j^T)\}_{j=1}^{N_T}$ where
$N_T \ll N_S$, our goal is to learn a classifier
$f_\theta: \mathcal{X} \rightarrow \mathcal{Y}$ that performs well
on both languages.

\section{Proposed Approach}

\subsection{Adaptive Source Selection}

We propose an adaptive method that selects source languages based on
typological similarity to the target language.  Given a set of
candidate source languages $\{S_1, \ldots, S_K\}$, we compute a
similarity score:

\begin{equation}
    \text{sim}(S_k, T) = \sum_{i=1}^{F} w_i \cdot \delta(t_i^{S_k}, t_i^T)
    \label{eq:similarity}
\end{equation}

where $t_i^L$ is the value of typological feature $i$ for language
$L$, $\delta$ is a distance function, and $w_i$ is the feature weight
learned from validation data.

\subsection{Fine-Tuning Strategy}

For each selected source language $S_k$, we fine-tune the pre-trained
model on $\mathcal{D}_{S_k}$ and then adapt to $\mathcal{D}_T$ using
elastic weight consolidation (EWC) to prevent catastrophic forgetting:

\begin{align}
    \mathcal{L}_{\text{EWC}}(\theta) &= \mathcal{L}_{\text{task}}(\theta; \mathcal{D}_T) \nonumber \\
    &\quad + \frac{\lambda}{2} \sum_{i} F_i (\theta_i - \theta_i^{*})^2
    \label{eq:ewc}
\end{align}

where $F_i$ is the Fisher information matrix diagonal and
$\theta^*$ are the fine-tuned parameters.

\subsection{Ensemble Combination}

The final prediction combines outputs from multiple source-language
models:

\begin{equation}
    \hat{y} = \arg\max_y \sum_{k=1}^{K} \alpha_k \cdot P(y | x^T; \theta_k)
    \label{eq:ensemble}
\end{equation}

where $\alpha_k$ are learned combination weights.

\section{Evaluation Framework}

We evaluate on three tasks: named entity recognition (NER), part-of-speech
tagging (POS), and sentiment analysis.  The evaluation includes 15
languages from 8 language families:

\begin{table}[htbp]
    \centering
    \caption{Languages and families in the evaluation set.}
    \label{tab:languages}
    \begin{tabular}{@{}llc@{}}
        \toprule
        Family & Language & Speakers (M) \\
        \midrule
        Indo-European & English, German, French & 1500 \\
        Indo-European & Russian, Hindi, Bengali & 900 \\
        Sino-Tibetan & Chinese, Burmese & 1200 \\
        Japonic & Japanese & 125 \\
        Koreanic & Korean & 80 \\
        Turkic & Turkish, Uzbek & 100 \\
        Niger-Congo & Swahili & 100 \\
        Austronesian & Indonesian, Tagalog & 300 \\
        \bottomrule
    \end{tabular}
\end{table}

\foreignlanguage{german}{Wir evaluieren auf drei Aufgaben: NER,
POS-Tagging und Sentiment-Analyse.  Die Evaluierung umfasst 15
Sprachen aus 8 Sprachfamilien.}
